\section{Introduction}\label{introduction}

Simulated 3D environments serve games, film, robot learning and embodied AI \citep{yang2024holodeck}. Across these uses they share a common structure of separate objects with collision geometry and mass, a ground with its gravity, light from identifiable sources, and scenes that can be reopened, modified and simulated again. Authoring such an environment by hand is expensive, and a generative system that takes over this work must leave the result in the hands of its author, since in the Game Developers Conference's 2026 survey 52 percent of game workers considered generative AI to be harming the industry \citep{gdc2026survey}.

World models generate an environment as a sequence of video frames \citep{bruce2024genie, genie3} that no physics simulator can step forward in time. 3D environment generators \citep{wang2025worldgen, hyworld2} produce geometry, mostly as a single fused asset from which separate objects have to be recovered afterwards \citep{hu2026worldact, niu2026worldsculpt}. Procedural systems \citep{raistrick2023infinigen, raistrick2024indoors} and coding agents that write scenes as executable programs \citep{hu2024scenecraft, li2026rcwm, guo2026worldclaw, yin2026autoue, he2026thinking, wang2026codeasworlds} keep the environment explicit and inspectable, but the models they produce often float above their supports or come apart \citep{gao2026codebench3d}, and such agents assess their own geometry unreliably \citep{li2026lego}. Open vision foundation models now recover segmentation masks \citep{carion2025sam3} and 3D shape \citep{li2026pixal3d, xiang2024trellis, trellis2} from single images, and a coding agent can call them as tools, yet none of them provides physical properties, and objects made by different models rarely share one visual style \citep{wu2026production, yang2026flowscene}.

The components needed for the physics already exist. OpenUSD \citep{openusd}, the scene description that Pixar developed for film production, composes scenes from smaller ones by reference and describes rigid bodies, colliders and mass in the same stage, and MuJoCo \citep{todorov2012mujoco} is a widely used engine for bodies in contact. What is missing is a route from what a creator intends, typically a single concept picture, to a complete environment in this form that still reads as that picture.

SCORE is a framework for that route. Its environments are complete, since every object, surface and light a scene shows exists wherever the viewer looks; owned, because every model in the framework permits commercial use of its output and every object comes from the creator's own asset library, score\_lib; and responsive, because every object carries its collider, mass and material, so a simulator steps the environment as written. This paper reports a change of method. SCORE's first method derived every scene from its concept through a chain of model readings, written plans and generated shells. Measured on 26 places of a first environment, that derivation produced scenes far from their concepts (Section 2). We replaced it with kit placement (Section 3): an agent lays library pieces, renders the concept's own view, compares and corrects, under a physics-readiness check that refuses any scene a simulator could not step. People are then routed through the checked floor and given musculoskeletal sim bodies (Section 4). Section 5 reports four places built this way, with their scores against the concepts, their checks, their walkers and their costs, and Section 6 states what still fails.
